\documentclass[10pt,a4paper]{amsart}
\usepackage[margin=19mm]{geometry}
\usepackage{amsmath,amssymb,mathtools}
\usepackage[scr=rsfs,cal=euler]{mathalfa}
\usepackage{xfrac}
\usepackage[unicode,hidelinks,bookmarks=false]{hyperref}
\newcommand{\ie}{i.e.\ }
\newcommand{\cf}{cf.\ }
\newcommand{\resp}{respectively\ }
\newcommand{\Subsection}{\subsection}
\providecommand{\nobreakdash}{\nobreak\hbox{-}}
\setlength{\emergencystretch}{2em}
\multlinegap0pt
\usepackage{amssymb}
\usepackage{mathtools}
\newtheorem{theorem}{Theorem}
\title{On $R$-equivalence of Automorphism groups of Associative algebras}
\author{Dibyendu Das}
\date{Source: arXiv:2512.24357v1 (CC BY 4.0)}
\begin{document}
\maketitle

\noindent\emph{Source and licence.} On $R$-equivalence of Automorphism groups of Associative algebras. Version arXiv:2512.24357v1 (CC BY 4.0), distributed by arXiv under the Creative Commons Attribution 4.0 International licence, http://creativecommons.org/licenses/by/4.0/, as recorded in the arXiv licence metadata for this exact version and shown on the abstract page https://arxiv.org/abs/2512.24357v1. Full licence notice: \texttt{licenses/arxiv-2512.24357-CC-BY-4.0.txt}.\par\smallskip

\noindent\textit{Editorial scope.} Counted unit: the article's theorem 3.1 (source0.tex lines 315--317). The article's complete original proof of it is delivered with it (source0.tex lines 329--331, one \texttt{proof} environment of 3 lines, which the article places away from the statement). No mathematics is written, repaired or re-derived: the statement and the proof are the article's own lines, in the article's own notation, and no retained line is substituted or re-flowed.  No further source-local context is needed: every cross-reference of every retained line resolves inside the two delivered spans.  Nothing was omitted from any retained span, so the package carries no omission record.  No retained line contains a citation, so the wrapper carries no bibliography and no bibliography entry.  No retained line contains a figure environment, an image-inclusion command or drawing code, so no image is delivered and the package declares no asset dependency.  Editorial formatting only, in addition: an AMS \texttt{amsart} preamble that restates the article's own declarations and macros used by the retained lines, namely the environments \texttt{theorem} on separate counters, together with the standard packages those lines need.  The retained environments print in this document as Theorem 1 in order of appearance, whereas the article prints the same items with the numbering of its own full document.  The complete native source of the article as distributed in the arXiv e-print for this exact version is bundled unchanged as \texttt{sources/191886/source0.tex}.
\begin{theorem}\label{Theorem 3.1}
Let $A$ be a split finite-dimensional associative algebra over a perfect field $k$, and let $J$ be the radical of $A$. If $U$ is the $k$-unipotent radical of $G_A$, then $J^2=0$ if and only if $J\subset A^{U}$, where $A^U\coloneq \{a\in A:\sigma(a)=a \quad \forall \sigma\in U(F), \text{ for all field extension $F$ over $k$ } \}$. 
\end{theorem}

\begin{proof}
    Let $G_A$ be reductive, then $U=\{Id\}$, $U$ is the $k$-unipotent radical. Hence, it follows that $A^U=A\supset J$. Therefore, by Theorem \ref{Theorem 3.1}, we have $J^2=0$. Since $U=\{Id\}$, we also have $\hat{J}=\{Id\}$, which further implies $J\subset Z(A)$. Conversely, if we assume $J^2=0$, then it follows from the proof of Theorem \ref{Theorem 3.1} that the $k$-unipotent radical of $G_A$ is $\hat{J}=U$. However, $J\subset Z(A)$ shows that $\hat{J}=\{Id\}$. Hence, the group $G_A$ is reductive.
\end{proof}

\end{document}
